\documentclass[11pt,a4paper]{article}
\usepackage[margin=2.3cm]{geometry}
\usepackage{amsmath,amssymb,booktabs,longtable,xcolor,fancyvrb}
\usepackage[colorlinks=true,linkcolor=blue!45!black,urlcolor=blue!45!black]{hyperref}
\usepackage{titlesec}
\titleformat{\section}{\large\bfseries}{\thesection}{0.6em}{}
\setlength{\parskip}{0.35em}
\newcommand{\file}[1]{\texttt{\small #1}}
\title{\bf Iteration 004: complete process record\\[3pt]
\large From simulated genotypes to the final analysis --- every step, in order}
\author{fmbenchmark}\date{\today}
\begin{document}\maketitle

\begin{abstract}\noindent
This document records exactly how the Iteration 004 fine-mapping benchmark
results were produced. It is written to be followed by someone who has not seen
the project before: every stage names the file that implements it, the
parameters it was run with, what it consumed and what it produced. Where a step
went wrong during execution and had to be redone, that is recorded too, in
\S\ref{sec:incidents} --- a process record that omits its failures cannot be
audited. The companion document \file{iter004\_full\_analysis.pdf} presents the
results; this one explains how they came to exist.
\end{abstract}

\tableofcontents\bigskip

\section{Overview}

The pipeline has nine stages. Each is a separate script; each writes its output
to disk before the next begins, so any stage can be re-run without repeating the
ones before it.

\begin{center}\small
\begin{tabular}{rlll}\toprule
& stage & implemented in & produces\\\midrule
1 & design grid & \file{generate\_params\_grid.R} & \file{params\_grid.csv}\\
2 & genotypes & \file{simulate\_genotypes.R} & region panel (in \file{sim.rds})\\
3 & annotations & \file{simulate\_phenotypes.R} & annotation matrices\\
4 & causal selection & \file{simulate\_phenotypes.R} & \file{truth\$causal\_indices}\\
5 & phenotypes & \file{simulate\_phenotypes.R} & $y$, $\beta$, realised PVE\\
6 & summary stats + LD & \file{simulate\_phenotypes.R} & $z$, LD matrix\\
7 & method fitting & \file{run\_methods.R} + 18 wrappers & \file{results.rds}\\
8 & collection & \file{iter004\_collect.R} & L1 fit table\\
9 & analysis & \file{iter004\_report.R} + 3 sense scripts & L2, L3, sense outputs\\
\bottomrule\end{tabular}\end{center}

\paragraph{Determinism.} Stages 2--6 run under \texttt{seed = 1000 + row\_index},
set once per design row. Given the same grid and the same VCF inputs, the
simulated data are reproducible exactly. Stage 9's cross-fitted test uses
\texttt{set.seed(20260806)}. Stage 7 is not seeded centrally: each method
controls its own randomness, and the BEATRICE family's Python trainer sets
\texttt{torch.manual\_seed(1)} internally.

\section{Software environment}

\begin{center}\small
\begin{tabular}{ll}\toprule
R & 4.5.2 (\texttt{R/4.5.2-gfbf-2025b})\\
Python & 3.12.3 (\texttt{Python/3.12.3-GCCcore-13.3.0})\\
GSL & 2.8 (\texttt{GSL/2.8-GCC-14.3.0})\\
scheduler & PBS Pro, queue \texttt{v1\_small72a}, 72\,h walltime\\
toolchain root & \file{/rds/general/project/neurogenomics-lab/live/tools}\\
\bottomrule\end{tabular}\end{center}

The toolchain root holds a Python virtualenv (\file{fmpy-venv}), a wrapper
(\file{py-venv-runner.sh}) that every Python-backed method is handed in place of
a bare interpreter, the compiled binaries \file{finemap\_v1.4.2} and
\file{PAINTOR}, and the source checkouts \file{SparsePro}, \file{fine-mapping-inf}
and \file{Funmap}. The last of these must be \emph{built}, not merely cloned: its
\file{finemapinf} sub-package declares a C extension.

\paragraph{This location matters.} It must be on permanent storage. An earlier
execution placed it behind a symlink onto a filesystem subject to a rolling
30-day purge; see \S\ref{sec:incidents}.

Rebuild it with:
\begin{Verbatim}[fontsize=\small,frame=single]
DEST=/rds/general/project/neurogenomics-lab/live/tools \
  bash scripts/hpc/rebuild_toolchain.sh
bash scripts/hpc/check_toolchain.sh      # must print RESULT: PASS
\end{Verbatim}

\file{check\_toolchain.sh} derives the required Python imports by walking the
tool sources with \texttt{ast}, subtracting the standard library, and importing
what remains. It does not use a hand-written package list, because a
hand-written list is what caused the failure in \S\ref{sec:incidents}.

\section{Stage 1 --- the design grid}

\file{scripts/hpc/generate\_params\_grid.R} writes 45 rows to
\file{params\_grid.csv}. Each row is one combination of architecture,
annotation type, enrichment fold and (for the infinitesimal arm)
$p_{\text{causal}}$. The within-row sweep is fixed:

\begin{center}\small
\begin{tabular}{ll}\toprule
$S$ & $1,2,3,5,10$\\
$\phi$ & $0.05, 0.1, 0.2, 0.4, 0.6$\\
iterations & 10\\
regions & 10, nominal sizes $\{500,500,750,750,1000,1000,1500,1500,2000,2000\}$\\
$n$ & 5000\\
LD & in-sample only (\texttt{n\_ref = NA})\\
annotations & 10 per region\\
\bottomrule\end{tabular}\end{center}

$45\times5\times5\times10 = 11{,}250$ scenarios; $\times10$ regions
$=112{,}500$ fits per method.

\paragraph{The grid is regenerated at every submission, never reused from disk.}
It is untracked and persists between iterations, and an earlier run silently
reused a stale 135-row grid from Iteration 003. The submitter now regenerates it
unconditionally and echoes the resulting design before calling \texttt{qsub}, so
a mismatch is caught in seconds rather than after 72 hours.

\section{Stage 2 --- genotypes}

\file{R/simulate\_genotypes.R} draws each region from real 1000 Genomes loci
under \file{data/vcf}, filtered to $\mathrm{MAF}\ge0.01$, then standardises
columns.

\paragraph{Regions can be smaller than requested.} If a locus does not contain
enough variants passing the filter, the region is whatever is available. Measured
over the 709 distinct regions actually used:

\begin{center}\small
\begin{tabular}{lrrr}\toprule
nominal & median actual & ratio & within 5\%\\\midrule
500 & 499 & 1.00 & 100\%\\
750 & 749 & 1.00 & 96\%\\
1000 & 999 & 1.00 & 97\%\\
1500 & 1483 & 0.99 & 62\%\\
2000 & 1590 & 0.80 & 26\%\\
\bottomrule\end{tabular}\end{center}

The analysis uses the \emph{nominal} size as its design factor. Variation in
actual size within a nominal class is absorbed into the between-region variance
component $\sigma^2_u$ and correctly subtracted by the noise correction, so the
design remains valid; only the label is approximate.

\paragraph{Regions are shared across the whole row.} All 250 scenarios of a
design row use the same ten regions, simulated once and cached in \file{sim.rds}.
This is what creates the split-plot structure the analysis corrects for, and it
is why two regions exist per size class --- without a replicate draw, $\sigma^2_u$
would not be identifiable.

\section{Stage 3 --- annotations}

Ten annotations per region. Binary annotations are Bernoulli indicators;
continuous ones are drawn from a shared-factor model so that annotations within
an enrichment group are correlated. The same matrix is passed to every
annotation-aware method, so no method sees information another does not.

\section{Stage 4 --- causal variant selection}

\file{select\_causal\_variants()} in \file{R/simulate\_phenotypes.R}.

\begin{itemize}
\item \textbf{No annotations}: $S$ indices drawn uniformly without replacement.
\item \textbf{With annotations}: weights
$w_j=\exp\!\big(A_j^\top\log\gamma\big)$, normalised to probabilities, then $S$
drawn without replacement.
\end{itemize}

$\gamma$ is the enrichment vector for that row. \file{polyfun\_oracle}
reconstructs $\pi_j\propto\exp(A_j^\top\log\gamma)$ from the stored truth,
using the identical formula --- which is what makes it a true ceiling rather
than a strong competitor.

\section{Stage 5 --- phenotypes}

\subsection{Sparse}
Effects $\beta\sim N(0,\sigma^2_\beta)$ on the $S$ causal variants; the residual
variance is then chosen to hit the target heritability exactly:
\[
\sigma^2_\varepsilon=\operatorname{Var}(X\beta)\frac{1-\phi}{\phi}
\quad\Longrightarrow\quad \mathrm{PVE}=\phi\ \text{exactly}.
\]

\subsection{Sparse plus infinitesimal}
A second component is drawn on the \textbf{non-causal variants only},
$\alpha_{nc}\sim N(0,1/m_{nc})$ (the BEATRICE formulation;
\texttt{inf\_model = "beatrice"}, the default, is what was used). The components
are then rescaled so that
\[
\operatorname{Var}(g_{\text{sparse}})=p_{\text{causal}}\phi,\quad
\operatorname{Var}(g_{\text{inf}})=(1-p_{\text{causal}})\phi,\quad
\sigma^2_\varepsilon=1-\phi,
\]
and $\beta$ is rescaled by the same factor so the recorded truth stays
consistent with the realised phenotype.

\paragraph{Consequence for interpretation.} Every non-causal variant then has a
genuine non-zero effect, while the ground truth used for scoring is the $S$
sparse causals only. A method that detects infinitesimal signal is counted as
making a false discovery. Absolute FDR is therefore not comparable between the
two arms; \emph{degradation} between arms is the meaningful quantity.

\section{Stage 6 --- summary statistics and LD}

\file{compute\_summary\_statistics()}: marginal OLS per variant on the centred
phenotype,
\[
\hat\beta_j=\frac{x_j^\top y_c}{x_j^\top x_j},\quad
\mathrm{RSS}_j=\|y_c\|^2-\hat\beta_j x_j^\top y_c,\quad
\mathrm{se}_j=\sqrt{\frac{\mathrm{RSS}_j}{(n-2)\,x_j^\top x_j}},\quad
z_j=\hat\beta_j/\mathrm{se}_j .
\]

LD is \texttt{cor(X)} on the \emph{same} genotype matrix that generated the
phenotype --- in-sample throughout, with no reference panel anywhere in this
iteration. Iteration 003 established that in-sample and reference-panel LD can
invert method rankings, so no conclusion here transfers to the panel setting.
\section{Stage 7 --- fitting the methods}

\file{scripts/hpc/run\_benchmark\_job.R} is the array worker. Task $t$ maps to
\[
\text{row}=\Big\lfloor\frac{t-1}{\text{tasks per row}}\Big\rfloor+1,
\qquad \text{tasks per row}=\Big\lceil\frac{250}{\text{chunk}}\Big\rceil ,
\]
with chunk $=5$, so 50 tasks per row and 2{,}250 tasks for the full grid. Each
task loads its row's cached \file{sim.rds}, then fits all methods to each of its
five scenarios in turn, writing \file{results.rds} per scenario before moving on.

\paragraph{Eighteen methods were fitted.}
\texttt{susie}, \texttt{susie\_inf}, \texttt{finemap}, \texttt{finemap\_inf},
\texttt{abf}, \texttt{marginal\_z}, \texttt{finimom}, \texttt{sparsepro},
\texttt{funmap}, \texttt{paintor}, \texttt{polyfun\_oracle},
\texttt{polyfun\_est}, \texttt{polyfun\_ldsc}, \texttt{sbayesrc},
\texttt{beatrice}, \texttt{functional\_beatrice}, \texttt{fb\_pooled},
\texttt{fb\_xregion}.

\paragraph{Four of these are reimplementations, not the upstream software.}
\texttt{polyfun\_oracle}, \texttt{polyfun\_est}, \texttt{polyfun\_ldsc} and
\texttt{sbayesrc} are implemented inside this package, because the published
tools depend on genome-wide reference infrastructure --- pre-computed UKB
annotation matrices, reference LD scores, eigen-decomposed LD folders --- that
does not exist for simulated per-locus data.

\begin{itemize}
\item \texttt{polyfun\_oracle} skips PolyFun's Stage 1 entirely and passes the
simulator's true prior to \texttt{susie\_rss()}. It is a ceiling, not PolyFun.
\item \texttt{polyfun\_est} is ``LDSC-lite'': $\chi^2$ regressed on each
variant's own annotations.
\item \texttt{polyfun\_ldsc} follows the canonical S-LDSC model, regressing on
annotation LD scores by weighted non-negative least squares. Closest of the three.
\item \texttt{sbayesrc} reimplements the SBayesRC \emph{algorithm} (spike plus
four normal slabs with annotation-modulated weights) in R.
\end{itemize}

Results for these four are evidence about the \emph{modelling strategy}, not
about the software of the same name. One difference is material enough to
affect interpretation: \texttt{sbayesrc}'s PIP is the fraction of MCMC samples
in which a variant occupies any non-spike component, and its smallest slab has
variance $5\times10^{-5}$ --- so a variant with an effect of no consequence
counts fully toward its PIP. That quantity answers ``does this variant have a
non-zero effect'', whereas every metric here scores against ``is this variant
one of the $S$ sparse causals''. Its mass ratios ($10$--$18$) and top-band
calibration ($0.996$ claimed, $0.621$ realised) follow from that mismatch of
definition rather than from any error, and are not comparable with the other
methods'.

CARMA was dropped part-way through on measured cost: $19{,}000$\,s per scenario
against $9{,}765$\,s for the other eighteen combined, i.e.\ $66\%$ of total
runtime, for a differentiator (LD-discrepancy outlier detection) that an
all-in-sample design cannot exercise. Scenarios fitted before that change do
contain CARMA results; they are retained on disk but excluded from every
analysis, because partial coverage of a balanced design unbalances every
stratum it touches.

\paragraph{The four BEATRICE variants.} All share one Binary Concrete
(Gumbel-Softmax) variational posterior over causal configurations and differ
only in where the prior comes from:

\begin{center}\small
\begin{tabular}{lll}\toprule
method & prior & fitted\\\midrule
\texttt{beatrice} & uniform & ---\\
\texttt{functional\_beatrice} & LassoNet on annotations & per region\\
\texttt{fb\_pooled} & shared logistic head & jointly, all regions\\
\texttt{fb\_xregion} & shared LassoNet head & jointly, all regions\\
\bottomrule\end{tabular}\end{center}

The two joint variants optimise
$L(\phi,\{\psi_r\})=\sum_r\mathrm{ELBO}_r(\psi_r;p_{0,r}=f_\phi(v_r))$ with a
\emph{single} optimiser over the shared head and all region posteriors: one
backward pass per step on $\tfrac1R\sum_r\mathrm{loss}_r$, so every parameter
moves simultaneously. It is not alternating optimisation. The equal weighting of
regions follows from the $1/R$ rather than being imposed separately.

With no annotations all four reduce to the same code path, which is why they are
bit-identical on both unannotated arms --- a mechanical negative control, not a
coincidence.

\paragraph{Hyperparameters} for the BEATRICE family come from Optuna trial
\#401: \texttt{sigma\_sq}~$=0.0899$, \texttt{sparse\_concrete}~$=51$,
\texttt{max\_iter}~$=2000$, \texttt{n\_caus}~$=10$, and for the functional
variants \texttt{prior\_regularisation}~$=0.1255$,
\texttt{lambda\_l1}~$=0.1223$, \texttt{hierarchy\_M}~$=10.15$. All competitors
run at their documented defaults. A separate analysis found the family's
sensitivity to these values to be small ($\pm0.004$ AP across a deliberately
wide hand-tuned sweep) relative to the architectural contrasts.

\paragraph{Per-task preflight.} Before any fitting, each task imports every
module the tool sources actually require --- the list derived by \texttt{ast}
walk, not written by hand --- and aborts if any fails. All 1{,}800 tasks of the
final run reported \texttt{preflight ok}.

\section{Stage 8 --- collection}

\file{scripts/analysis/iter004\_collect.R}, one array task per design row,
re-derives every metric from the raw PIPs in \file{results.rds} and the truth in
\file{sim.rds}. It does not read the \file{evaluation.rds} written during
fitting.

Two details matter for correctness:

\begin{itemize}
\item \textbf{The scenario index comes from the directory name}, not from
\texttt{fit\$scenario\_id}. The worker subsets \file{sim} to one scenario and
resets that field to 1, because \texttt{evaluate\_methods()} uses it as a list
index; the true index survives only in the \file{scenario\_<n>} path. Indexing
by the stored field would silently attach every fit to scenario 1.
\item \textbf{A fit with any non-finite PIP is marked failed, not repaired.} Its
total mass and the ranking of the affected variants are undefined, so no metric
derived from it is trustworthy. Imputing to zero would invent posterior mass and
flatter precisely the method that failed. A separate \texttt{bad\_pip} column,
scoped to fits that did \emph{not} declare an error, makes silent failure
countable and distinguishable from an honest one.
\end{itemize}

Output: the L1 table, one row per (scenario, region, method), $2{,}031{,}230$
rows, carrying $\sim45$ metric columns.

\section{Stage 9 --- aggregation and analysis}

\file{scripts/analysis/iter004\_report.R} binds the L1 partials, then:

\begin{enumerate}
\item Joins the grid, and \textbf{recodes \texttt{NA}-as-level before any
grouping}. \texttt{p\_causal} is \texttt{NA} on every sparse row and
\texttt{enrichment\_fold} is \texttt{NA} on every unannotated row;
\texttt{aggregate()}, \texttt{split()} and \texttt{table()} all drop \texttt{NA}
groups silently, which would delete those arms entirely.
\item Builds L2 (replicate level: cell $\times$ iteration $\times$ method) and
L3 (cell means), using a hand-rolled grouped reduction. \texttt{aggregate()} did
not complete in 120\,s at this scale; the replacement runs in 2.7\,s and was
verified to give identical results.
\item Runs the validity checks of \S\ref{sec:checks}.
\item Drives the three sense scripts.
\end{enumerate}

The three senses are defined and derived in full in the companion analysis
document; in outline, Sense V is a noise-corrected variance decomposition on
cell means, Sense D a cross-fitted test of whether a cell's winner is real, and
Sense G a set of paired within-replicate contrasts.

\section{Validity checks performed}\label{sec:checks}

\subsection{Automatic, on every run}
\begin{enumerate}
\item In-sample LD only --- zero rows with a reference-panel size.
\item 45 design rows present, matching the grid.
\item Balanced cells per analysed method --- $5{,}625$ for all 18.
\item Exclusions (CARMA) reported explicitly rather than silently applied.
\item No method entirely absent outside the unannotated arms.
\item \texttt{polyfun\_est} $=$ \texttt{susie} on the unannotated arms.
\item \texttt{polyfun\_ldsc} $=$ \texttt{susie} on the unannotated arms.
\end{enumerate}

Checks 6 and 7 are the strongest structural tests available: both methods reduce
to SuSiE analytically without annotations, so a discrepancy would indicate a
fault in the annotation plumbing rather than in a method.

\subsection{Additional checks performed for this report}
\begin{enumerate}\setcounter{enumi}{7}
\item Truth consistency: $\texttt{n\_causal}=S$ for all $2{,}031{,}230$ fit
rows. Zero mismatches.
\item Grid completeness: exactly $11{,}250$ distinct
$(\text{row},S,\phi,\text{iter})$ tuples; \texttt{region\_idx} $\in\{1,2\}$ only.
\item Metric internal consistency, per fit: $\mathrm{tp}_t\le n_t$ and
$\mathrm{tp}_t\le S$ at $t\in\{0.5,0.9,0.99\}$; $n_t$ monotone decreasing in
$t$; $\mathrm{AP}\in[0,1]$; $|CS_{95}|\le p$; PIP-band counts summing to $p$;
causal band counts summing to $S$. Zero violations of any.
\item Negative controls, exactly: all four BEATRICE variants identical to
machine precision on both unannotated arms; \texttt{polyfun\_\{est,ldsc,oracle\}}
identical to \texttt{susie} there.
\item Directional sanity: AP strictly decreasing in $S$; strictly increasing in
$p_{\text{causal}}$; sparse arm easier than sparse+infinitesimal. AP is
\emph{not} monotone in $\phi$ --- investigated and found to be a genuine
method-specific effect, not a fault.
\item Simulation code review: PVE calibration exact by construction in both
architectures; oracle prior uses the identical formula to causal selection.
\item Headroom validity: the oracle's ranking advantage over the baseline is
intact in all four annotated strata, so $\kappa$ is defined there.
\end{enumerate}

\section{Incidents during execution}\label{sec:incidents}

Recorded because a process record that omits its failures cannot be audited, and
because each fix is now a guard that a future run inherits.

\paragraph{1. Stale grid silently reused.} \file{params\_grid.csv} is untracked
and persists between iterations; a submission reused Iteration 003's 135-row
grid. \emph{Fix:} the submitter regenerates it unconditionally and prints the
design before submitting.

\paragraph{2. CARMA dominating runtime.} $66\%$ of total compute for an
untestable capability. \emph{Fix:} removed from the method set; prior partial
results retained but excluded from analysis.

\paragraph{3. Toolchain on a purging filesystem.} \file{\textasciitilde/tools}
was a symlink onto ephemeral storage subject to a rolling 30-day purge. It
decayed \emph{during} a run: activating the virtualenv kept succeeding while
imports inside it stopped resolving, so seven Python-backed methods recorded
$38$--$50\%$ failed fits for hours without the job stopping. \emph{Fix:}
toolchain moved to permanent project storage; a per-task preflight now aborts
immediately if any required import fails. Affected results were discarded and
the arm refitted.

\paragraph{4. Dependency list written from memory.} The rebuilt environment
installed five packages; BEATRICE also requires \texttt{imageio},
\texttt{seaborn} and \texttt{tqdm}, and FINEMAP-inf requires \texttt{bgzip}. All
fits died on \texttt{ModuleNotFoundError} after a check that had passed.
\emph{Fix:} the required imports are now derived from the tool sources by
\texttt{ast} walk and each is actually imported, not merely located --- an
installed-but-broken package passes \texttt{find\_spec} and fails at runtime.

\paragraph{5. \texttt{setuptools} 81 removed \texttt{pkg\_resources}.} An
unpinned upgrade broke FINEMAP-inf, which imports it. \emph{Fix:} pinned
\texttt{setuptools<81}. Separately, \file{fine-mapping-inf} must be built rather
than cloned (it declares a C extension) and its build needs
\texttt{--no-build-isolation} because \file{setup.py} imports numpy at build
time.

\paragraph{6. Silent all-NA output.} \texttt{finemap\_inf} returns an all-NA
posterior without raising an error at $\phi=0.6$, $S=1$ in the sparse arm. The
old evaluator crashed on it; the collect stage treated it as a successful fit.
\emph{Fix:} non-finite PIPs mark the fit failed and set a \texttt{bad\_pip} flag,
scoped to fits that declared no error, so silent failure is countable.

\paragraph{7. Scheduler instability.} The final tail of the array was
system-held twice for repeated launch failures, during a period when the site
had disabled \texttt{qstat} for slowness. \emph{Fix (operational):} the held
subjobs were deleted and resubmitted as fresh jobs; all completed. No data was
affected --- completed scenarios were already on disk.

\section{Reproducing the run}

From the project root on the cluster:

\begin{Verbatim}[fontsize=\small,frame=single]
# 0. toolchain (once)
DEST=/rds/general/project/neurogenomics-lab/live/tools \
  bash scripts/hpc/rebuild_toolchain.sh
bash scripts/hpc/check_toolchain.sh          # must print RESULT: PASS

# 1-7. simulate and fit  (2250 tasks, ~72h with wide concurrency)
module load R/4.5.2-gfbf-2025b
FMB_SCENARIOS_PER_TASK=5 PBS_QUEUE=v1_small72a PBS_WALLTIME=72:00:00 \
  FMB_SCRATCH=$EPHEMERAL/fmbench_iter004 \
  bash scripts/hpc/submit_benchmark_pbs.sh

# 8-9. verify, then collect and analyse  (submits itself only if the run is sound)
qsub -v PROJECT=$HOME/fine-mapping-benchmark \
  scripts/analysis/verify_and_analyse.sh
\end{Verbatim}

The verify step checks \file{results.rds} completeness against 11{,}250, scans
every task log for preflight failures, and gates on per-method declared-failure
and silent-NA rates before submitting the analysis. It refuses to proceed on a
short or damaged grid.

\section{Output manifest}

\begin{center}\small
\begin{tabular}{lll}\toprule
file & level & contents\\\midrule
\file{combined\_fit\_metrics.rds} & L1 & one row per (scenario, region, method)\\
\file{combined\_replicate\_metrics.rds} & L2 & one row per (cell, iteration, method)\\
\file{combined\_scenario\_metrics.rds} & L3 & cell means; the ANOVA input\\
\file{validity\_checks.txt} & --- & the checks of \S\ref{sec:checks}\\
\file{iter004\_sense\_v/} & --- & variance decomposition\\
\file{iter004\_sense\_d/} & --- & decision analysis, per-cell winners\\
\file{iter004\_sense\_g/} & --- & headroom, $\kappa$, paired contrasts\\
\bottomrule\end{tabular}\end{center}

L1 is the pivot: every statistic in the analysis document is re-derivable from
it without refitting. The raw PIPs it was built from live under
\file{fmbench\_iter004/results/benchmark} and should be archived off ephemeral
storage if any metric not already in L1 might be wanted later.

\end{document}
